\section{La inteligencia y el producto importan por igual}

Esta sección aborda la segunda palabra clave: producto. Durante los últimos tres años, la industria se volcó de forma casi obsesiva a explorar el límite superior de la inteligencia, pero ahora el producto vuelve a ser importante. ChatGPT tiene muchas barreras de entrada no técnicas, mientras que una herramienta de Coding o una empresa de modelos que solo cuenta con barreras técnicas puede ser alcanzada con facilidad. La fortaleza de OpenAI está en el equilibrio: por un lado explora el límite superior de la inteligencia y, por otro, convierte los beneficios de esa inteligencia en tráfico de producto y en presencia de marca en la mente de los usuarios.

\begin{figure}[H]
\centering
\includegraphics[width=0.96\textwidth]{intelligence-product-balance.png}
\caption{La inteligencia y el producto importan por igual: la exploración del límite superior de la inteligencia debe convertirse en tráfico, marca y presencia en la mente de los usuarios. Diagrama conceptual propio, elaborado a partir de la conversación de 00:33:35--00:38:52.}
\end{figure}

\begin{knowledgebox}{Lectura de la figura: la inteligencia no equivale automáticamente a un producto}
La capacidad del modelo aporta posibilidades; la experiencia de producto convierte esas posibilidades en uso predeterminado; el tráfico y la presencia de marca en la mente de los usuarios hacen que estos regresen de forma continua; la comercialización convierte la inteligencia en una capacidad de la empresa. Si falta cualquiera de estos eslabones, el resultado puede quedarse en una simple demo técnica.
\end{knowledgebox}

\subsection{Barreras del modelo y barreras del producto}

La sección anterior subrayó que el producto vuelve a ser importante; esta separa dos tipos de barreras. Las barreras del modelo provienen de los benchmark, el coding, las capacidades agentic, el contexto y el razonamiento; las barreras del producto provienen de los hábitos de los usuarios, la marca, los puntos de entrada, el retorno de datos y la comercialización. Las barreras del modelo pueden ser alcanzadas y las del producto también pueden ser absorbidas, pero solo la suma de ambas forma una empresa sólida.

\begin{figure}[H]
\centering
\includegraphics[width=0.94\textwidth]{model-vs-product-moat.png}
\caption{Barreras del modelo vs barreras del producto: las barreras técnicas pueden ser alcanzadas; las barreras del producto incluyen el tráfico, la marca y el punto de entrada predeterminado. Diagrama conceptual propio, elaborado a partir de la conversación de 00:33:35--00:38:52.}
\end{figure}

\begin{warningbox}{No hay que fijarse solo en las barreras técnicas}
Las barreras técnicas de un producto de IA pueden alcanzarse en un plazo de 3 a 6 meses. Si no hay retorno de datos, hábitos de los usuarios, canales, marca ni flujos de trabajo consolidados, la ventaja técnica puede convertirse fácilmente en una ventana pasajera.
\end{warningbox}

\begin{knowledgebox}{La parte no técnica de las barreras del producto}
\begin{tabularx}{\textwidth}{p{0.22\textwidth}p{0.34\textwidth}X}
\toprule
Barrera & Manifestación & Por qué es difícil de replicar \\
\midrule
Punto de entrada predeterminado & El usuario lo abre de forma natural todos los días & Requiere una marca y hábitos de uso construidos a largo plazo. \\
Integración en el flujo de trabajo & El producto entra en los procesos del equipo & El costo de sustitución es alto y los datos regresan de forma continua. \\
Presencia de marca en la mente del usuario & El usuario le asigna por omisión cierto tipo de tareas & El primer Magic Moment se recuerda durante mucho tiempo. \\
Comercialización & Cuenta con pagos y canales estables & No basta con un modelo potente; también hay que saber vender y presentar el producto. \\
\bottomrule
\end{tabularx}
\end{knowledgebox}

\subsection{Resumen de la sección}

La relación entre la inteligencia y el producto se ha vuelto más equilibrada. Las empresas de modelos necesitan la conversión en producto, y las empresas de producto también necesitan entender los modelos; ni la tecnología por sí sola ni el producto por sí solo son suficientes.
